% SPDX-FileCopyrightText: 2026 Will Cook
% SPDX-License-Identifier: CC-BY-4.0
\documentclass[11pt]{article}
\usepackage{paper-house-style}
\usepackage{booktabs,tabularx,array,xurl}
\usepackage[colorlinks=true,linkcolor=linkinternal,urlcolor=linkexternal,citecolor=linkinternal]{hyperref}
\hypersetup{pdftitle={Writing Mathematics from the Literature and Reviewed Revisions},pdfauthor={Will Cook},pdflang={en-GB},bookmarksnumbered=true}
\runninghead{Writing mathematics from reviewed revisions}
\title{Writing Mathematics from the Literature\\and Reviewed Revisions}
\subtitle{A method for reading, revising and retaining useful editorial lessons}
\author{Will Cook}
\date{30 September 2026}
\emergencystretch=1.8em
\newcommand{\caseid}[1]{\texttt{#1}}
\newcommand{\corpus}{\url{https://github.com/wcook04/plectis-erdos}}
\begin{document}
\maketitle
\paperprovenance
\begin{abstract}
A mathematical sentence can be easy to read and still conceal the step that
justifies it. Nearby papers offer more useful models than a list of preferred
words: they show how the field introduces hypotheses, names objects, connects
deductions and allocates space to a proof's difficult transition. We describe a
method for studying those choices and testing their use in another manuscript.
The examples come from reviewed revisions of eight Erd\H{o}s-problem papers.
They include an infinite construction whose remainder must vanish, an estimate
whose required accuracy should be explained before it is proved, and proposed
rules that became too strong when detached from their original case. Each lesson
retains its source, the old and proposed passages, the integrating decision and
a limit on reuse. The resulting guidance extends an existing public writing
skill and can be included at an exact version in a later review packet. This is
a documented editorial practice; it has not been evaluated as an intervention
in reader understanding.
\end{abstract}
\input{exposition/parts/reading.tex}
\input{exposition/parts/revisions.tex}
\input{exposition/parts/review.tex}
\input{exposition/parts/practice.tex}
\input{exposition/parts/references.tex}
\end{document}
